
\begin{table}[H]
   \caption{\label{tab:parto_adequado} Parto Adequado DiD (does NOT survive the parallel-trends test)}
   \centering
\small\setlength{\tabcolsep}{5pt}
   \begin{tabular}{lcc}
      \toprule
      Dependent Variable: & \multicolumn{2}{c}{rate}\\
                                                        & Private (SINASC, annual) & Private (TISS, quarterly) \\   
      Model:                                            & (1)                      & (2)\\  
      \midrule
      \emph{Variables}\\
      Post (2017+) $\times$ Treated municipality $=$ 1  & -0.0659$^{***}$          & -0.0159\\   
                                                        & (0.0146)                 & (0.0136)\\   
\addlinespace[2pt]
      \midrule
      \emph{Fixed-effects}\\
      Municipality                                      & Yes                      & Yes\\  
      Year                                              & Yes                      & \\  
      Quarter                                           &                          & Yes\\  
      \midrule
      \emph{Fit statistics}\\
      Observations                                      & 7,660                    & 22,303\\  
      R$^2$                                             & 0.856                    & 0.776\\  
\bottomrule
   \end{tabular}
   
   \par \raggedright 
   \footnotesize\textit{Notes:} Treated = municipality with a participating \emph{Parto Adequado} Fase 2 private hospital (a diluted exposure; TISS has no hospital identifier). Weighted by private births (SINASC) / deliveries (TISS). SE clustered by municipality. \emph{The SINASC event study rejects parallel pre-trends} (2010--2015 joint test $F=4.5$, $p<0.001$): treated munis are on a pre-existing differential downward trend, so this DiD is not interpreted causally. \footnotesize Significance: *** p$<$0.01, ** p$<$0.05, * p$<$0.10.
\end{table}


